% !TeX root = ../main.tex
\chapter{EPICS Base and IOCs}
\label{chap:epics_iocs}

\section{Introduction}
There are two IOCs associated with the calibration and test of the zPSC: The ATE IOC, and the zPSC IOC.

\section{EPICS Installation Overview}
	EPICS installations performed by M. Capotosto loosely follow the below directory structure:

	\begin{itemize}
		\item \texttt{EPICS\_BASE=/epics/base} or \texttt{/epics/base\_[version]} or similar
		\item \texttt{SUPPORT=/epics/support} -- or sometimes MODULES instead of SUPPORT is used.
			\par\bigskip
			\quad Modules required for the ATE and PSC are installed here:
			\begin{itemize}[topsep=0pt]
				\item seq: SNL-Sequencer
				\item pscdrv: Portable Streaming Controller Driver
				\item asyn: EPICS module for driver and device support
			\end{itemize}
		\item \texttt{epics/iocs} - each IOC is cloned into this repository
		\item \texttt{epics/systemd-softioc} - This is an NSLS-II utility for managing IOCs as a system service
	\end{itemize}

\section{Installing EPICS Base and Required Modules}
	\begin{enumerate}
		\item Install Dependencies for Build
			\hspace*{2em} Install git, build-essential, and perl:\\ \texttt{sudo apt -y install git build-essential perl libtirpc-dev re2c libevent-dev autoconf automake libtool pkg-config}
		\item Create the EPICS user group, and the empty EPICS directory: \\
			\hspace*{2em} \texttt{sudo groupadd epics}\\
			\hspace*{2em} \texttt{sudo usermod -aG epics \$USER}\\
			\hspace*{2em} \texttt{newgrp epics}\\
			\hspace*{2em} \texttt{sudo mkdir -p /epics}\\
			\hspace*{2em} \texttt{sudo chown \$USER:epics /epics}\\
			\hspace*{2em} \texttt{sudo chmod 2775 /epics}\\
		\item EPICS Base is cloned into \texttt{/epics/base} via:\\
			\hspace*{2em} \texttt{git clone --recursive -b 7.0} \\ \hspace*{4em}\texttt{https://github.com/epics-base/epics-base.git /epics/base}\\
			\hspace*{2em} \texttt{cd /epics/base}\\

            \textbf{Note:} The \texttt{7.0} branch tracks ongoing EPICS Base development and may contain commits newer than the most recent stable release. During testing, the development branch produced a Channel Access regression that caused \texttt{caput} operations to disconnect from the IOC. EPICS Base was therefore pinned to the stable \texttt{R7.0.10} release tag before building:

            \hspace*{2em} \texttt{git checkout R7.0.10}

            \hspace*{2em} \texttt{git submodule update --init --recursive}
		\item Build EPICS:\\
			\hspace*{2em} \texttt{make clean}\\
			\hspace*{2em} \texttt{make}\\

        \item Add the EPICS Base binaries to the user's PATH:\

            \hspace*{2em} \texttt{echo 'export EPICS\_BASE=/epics/base' >> \textasciitilde/.bashrc}\

            \hspace*{2em} \texttt{echo 'export EPICS\_HOST\_ARCH=linux-x86\_64' >> \textasciitilde/.bashrc}\

            \hspace*{2em} \texttt{echo 'export PATH=\$EPICS\_BASE/bin/\$EPICS\_HOST\_ARCH:\$PATH' >>    \textasciitilde/.bashrc}\

            \hspace*{2em} \texttt{source \textasciitilde/.bashrc}\

		\item Create the remaining directory structure:\\
			\hspace*{2em} \texttt{mkdir -p /epics/iocs /epics/modules}\\

		\item Clone asyn:
			\begin{enumerate}
				\item \texttt{cd /epics/modules}
				\item \texttt{git clone https://github.com/epics-modules/asyn.git asyn}
				\item \texttt{cd /epics/modules/asyn}
				\item Set up required files in \texttt{configure/RELEASE, configure/RELEASE.LOCAL}
				\item Build with \texttt{make clean}, and \texttt{make}
			\end{enumerate}

		\item Clone Sequencer:
		\begin{enumerate}
			\item \texttt{cd /epics/modules}
			\item \texttt{git clone https://github.com/epics-modules/sequencer.git seq}
			\item \texttt{cd /epics/modules/seq}
			\item Set up required files in \texttt{configure/RELEASE, configure/RELEASE.LOCAL}
			\item Build with \texttt{make clean}, and \texttt{make}
		\end{enumerate}

		\item Clone PSC Driver:
		\begin{enumerate}
			\item \texttt{cd /epics/modules}
			\item \texttt{git clone https://github.com/mdavidsaver/pscdrv.git pscdrv}
			\item \texttt{cd /epics/modules/pscdrv}
			\item Set up required files in \texttt{configure/RELEASE, configure/RELEASE.LOCAL}
			\item Build with \texttt{make clean}, and \texttt{make}
		\end{enumerate}

		\item Clone Manage IOCs Utility:
		\begin{enumerate}
			\item \texttt{cd /epics}
			\item \texttt{git clone https://github.com/NSLS2/systemd-softioc.git}
			\item Navigate to the repo directory \texttt{cd systemd-softioc}
			\item Make the install script executable and run:\\
			\hspace*{2em}\texttt{sudo chmod +x install.sh}\\
			\hspace*{2em}\texttt{sudo ./install.sh}\\
			\hspace*{2em}\texttt{sudo usermod -aG epics softioc}\\
			\hspace*{2em}\texttt{newgrp}\\
			\hspace*{2em}\texttt{sudo chown softioc:epics /epics/iocs}\\
			\hspace*{2em}\texttt{sudo chmod 2775 /epics/iocs}\\
		\end{enumerate}
	\end{enumerate}

\section{Installing the ATE IOC}
	\begin{enumerate}
		\item Change to IOC directory \texttt{cd /epics/iocs}
		\item Clone the repo:\\
		\hspace*{1em} \texttt{git clone https://github.com/capotosto/ALSu-PSC\_ATE-ioc.git PSCtest}
		\item \texttt{cd PSCtest}
		\item \texttt{git switch Released,-UPC/BPC-Split-IOC-Ver-2}
		\item Create the \texttt{RELEASE.local} file, \texttt{touch configure/RELEASE.local}
		\item \texttt{echo -e "EPICS\_BASE=/epics/base\textbackslash nASYN=/epics/modules/asyn" >}\\
			\hspace*{2em} \texttt{ /epics/iocs/PSCtest/configure/RELEASE.local}
		\item Build the IOC: \texttt{make clean}, \texttt{make}
		\item Update the st.cmd file: \texttt{cd iocBoot/iocPSCtest}
			\begin{enumerate}
			\item Edit the st.cmd file: \texttt{nano st.cmd}
			\item Update the EPICS\_BASE variable if needed.
			\item Update the IPPORT to match the ATE. Repository has 10.69.26.4:5000 as default. Changing this later from the Phoebus/CSS panels may not work!
		\end{enumerate}

		\item Find the next available port: \texttt{manage-iocs nextport}
		\item Configure this in the config file: \texttt{nano config}
		\item Set the \texttt{NAME=PSCtest}, \texttt{PORT=} your next port, \texttt{USER=softioc}, and \texttt{HOST=\$HOSTNAME}
		\item Update the directory path in the root-level st.cmd to point to the right binaries
		\item Enroll in manage-iocs as a service: \texttt{sudo manage-iocs install PSCtest}
		\item Enable auto-start on boot: \texttt{sudo manage-iocs enable PSCtest}
		\item Check the status, it should be running! \texttt{manage-iocs status}
		\item And if you want other information on the ioc: \texttt{manage-iocs report}
	\end{enumerate}

\section{Installing the PSC IOC}
\begin{enumerate}
	\item Change to IOC directory \texttt{cd /epics/iocs}
	\item Clone the repo:\\
	\hspace*{1em} \texttt{git clone https://github.com/jamead/zpsc-ioc.git zpsc\_ioc}
	\item \texttt{cd zpsc\_ioc}
	\item Create the \texttt{RELEASE.local} file, \texttt{touch configure/RELEASE.local}
	\item \texttt{echo -e "EPICS\_BASE = /epics/base\textbackslash n}\\
		\hspace*{2em}\texttt{PSCDRV = /epics/modules/pscdrv\textbackslash n}\\
		\hspace*{2em}\texttt{SEQ = /epics/modules/seq" >}\\
		\hspace*{2em} \texttt{/epics/iocs/zpsc\_ioc/configure/RELEASE.local}
	\item Build the IOC: \texttt{make clean}, \texttt{make}
	\item Update the st.cmd file as needed: \texttt{cd iocBoot/ioczpsc}
	\item Create the st.cmd file used by systemd-softioc: \begin{verbatim}echo -e '#!/bin/bash\n\n
		cd "/epics/iocs/zpsc_ioc/iocBoot/ioczpsc"\n./st.cmd' > st.cmd\end{verbatim}
	\item Make the new st.cmd executable: \texttt{chmod +x st.cmd}
	\item Find the next available port: \texttt{manage-iocs nextport}
	\item Configure this in the config file: \texttt{nano config}
	\item Set the \texttt{NAME=zpsc\_ioc}, \texttt{PORT=} your next port, \texttt{USER=softioc}, and \texttt{HOST=\$HOSTNAME}
	\item Update the directory path in the root-level st.cmd to the \texttt{iocBoot/ioczpsc/st.cmd}
	\item Enroll in manage-iocs as a service: \texttt{sudo manage-iocs install zpsc\_ioc}
	\item Enable auto-start on boot: \texttt{sudo manage-iocs enable zpsc\_ioc}
	\item Check the status, it should be running! \texttt{manage-iocs status}
	\item And if you want other information on the ioc: \texttt{manage-iocs report}
\end{enumerate}




















